\documentclass[
  % doi=10.XXXX/XXXXXX,   % Uncomment and fill in for the camera-ready version.
  % startpage=1,           % Uncomment and set first page number for the camera-ready version.
  year=26,                 % Two-digit year (FMCAD 20YY).
]{fmcad}

\usepackage[T1]{fontenc}
\usepackage[utf8]{inputenc}
\usepackage{microtype}
\usepackage{cite}
\usepackage{amsmath,amssymb,amsfonts}
\usepackage{algorithm}
\usepackage{algorithmic}
\usepackage{graphicx}
\usepackage{subcaption}
\usepackage{textcomp}
\usepackage{xcolor}
\usepackage{booktabs} 
\usepackage[square,comma,numbers,sort&compress]{natbib}      % Offers more control over citation appearance and bib spacing

% Relax float-placement parameters so figures don't get pushed to the end.
\renewcommand{\topfraction}{0.95}
\renewcommand{\bottomfraction}{0.95}
\renewcommand{\textfraction}{0.05}
\renewcommand{\floatpagefraction}{0.7}
\renewcommand{\dbltopfraction}{0.95}
\renewcommand{\dblfloatpagefraction}{0.7}
\setcounter{topnumber}{3}
\setcounter{bottomnumber}{3}
\setcounter{totalnumber}{6}
\setcounter{dbltopnumber}{3}

% --- todo notes ---
% Comment out the next line (or pass [disable]) before camera-ready
% to suppress every \todo / \zk / \as marker in one go.
\usepackage[colorinlistoftodos,prependcaption,textsize=small]{todonotes}
% \usepackage[disable]{todonotes}

% Per-author inline todo macros. Usage: \zk{rephrase this} or \as{cite?}.
\newcommand{\zk}[1]{\todo[inline,color=blue!20,size=\small]{\textbf{ZK:} #1}}
\newcommand{\as}[1]{\todo[inline,color=green!20,size=\small]{\textbf{AS:} #1}}
% Margin variants (shorter notes that don't break the paragraph).
\newcommand{\zkm}[1]{\todo[color=blue!20,size=\scriptsize]{ZK: #1}}
\newcommand{\asm}[1]{\todo[color=green!20,size=\scriptsize]{AS: #1}}


\begin{document}

\title{Parallelizing Congruence Closure}

\IEEEoverridecommandlockouts
\author{
  \IEEEauthorblockN{Zachary Kent\textsuperscript{*},
                    Amar Shah\textsuperscript{*}\thanks{\textsuperscript{*}Both authors contributed equally to this research.}}
  \IEEEauthorblockA{Carnegie Mellon University, Pittsburgh, PA, USA \\
                    \texttt{\{zkent, amarshah\}@cmu.edu}}
}

\maketitle

\begin{abstract}
  Given a set of ground equalities over terms with uninterpreted function
  symbols, congruence closure computes all ground equalities entailed by
  reflexivity, symmetry, transitivity, and congruence. While congruence closure
  is a performance-critical component of satisfiability modulo theories (SMT)
  solvers, automatic theorem provers, and equality saturation engines, all
  existing implementations are sequential. In this work, we provide the first
  parallel algorithm for congruence closure. 
  
  We build on two ideas from the parallel algorithms literature: a union-find
  data structure that supports concurrent union operations and a parallel
  semisort (group-by) algorithm that allows us to discover new congruences in
  parallel.

  We present two algorithms. The first replaces the classic sequential worklist
  with a bulk-synchronous closure loop: it processes all pending unions in
  parallel, then detects new congruences, and repeats until a fixpoint. The
  second removes the global barrier between these phases, detecting congruences
  concurrently with union processing and feeding the resulting unions back into
  a shared worklist. We evaluate both these algorithms against a sequential
  baseline on a number of synthetic and real-world equality saturation
  benchmarks. Both of our algorithms can be extended to support backtracking,
  proof production, and dynamic term insertion. These algorithms are compatible
  with modern SMT solver and equality saturation architectures.
\end{abstract}

\cite{uf}

\section{Introduction}
\label{sec:intro}

Congruence closure is an important algorithm for satisfiability modulo theories
(SMT) solving~\cite{smt}, equality saturation~\cite{eqsat}, type
unification~\cite{unification}, and hardware verification~\cite{clausal-cc}.
Congruence closure determines whether a set of equalities and disequalities
between terms is satisfiable. For example, given the equalities $a = b$ and
$b = c$ and the disequality $f(a) \neq f(c)$, congruence closure will
conclude that this is unsatisfiable: $a = c$ follows by \emph{transitivity},
and then $f(a) = f(c)$ follows by \emph{congruence}, contradicting the
disequality. To decide this problem in practice, one must maintain
equalities using a union-find data structure and propagate congruences to
parent terms.


% Satisfiability modulo theories (SMT) solvers are used to decided the
% satisfiability of first-order logic formulas. State-of-the-art solvers such as
% Z3~\cite{z3} and cvc5~\cite{cvc5} can reason about a wide variety of
% ``theories'' such as uninterpreted functions, linear arithmetic, bit-vectors,
% and many others.

Congruence closure is a component of SAT solvers like Kissat~\cite{kissat};
equality saturation engines like egg~\cite{egg} and egglog~\cite{egglog}; and
SMT solvers like Z3~\cite{z3} and cvc5~\cite{cvc5}. These tools are used in a
variety of applications, including circuit equivalence
checking~\cite{clausal-cc}, software verification~\cite{dafny,verus}, hardware
design~\cite{hardware}, and mathematical theorem proving~\cite{lean-grind}.
Hence, their performance is critical, and much work has gone into improving
their efficiency. One promising avenue for improving the performance is
\emph{parallelization}: utilizing multiple processors to solve the problem
faster. 

To this end, we present two novel parallel congruence closure algorithms. Our
algorithms use a concurrent union-find data structure and a bulk-synchronous
parallel design to efficiently compute the congruence closure of a set of
equalities and disequalities.

% Sequential congruence closure algorithms are based on the seminal work by Nelson
% and Oppen~\cite{nelson-oppen-cc}, which maintains a worklist of pairs of terms
% that are known to be equal. The algorithm repeatedly processes pairs from the
% worklist, merging their equivalence classes in the union-find data structure and
% adding new pairs to the worklist until no more pairs can be added.

At a high level, our algorithms replace the classical sequential worklist of
Nelson and Oppen~\cite{nelson-oppen-cc} with a \emph{bulk-synchronous} closure
loop. Initial equalities are unioned into a lock-free concurrent
union-find~\cite{uf}. The algorithm then proceeds in rounds: each round
identifies a frontier of terms whose congruence class may have changed,
discovers all new congruences among them in parallel, and unions the resulting
classes. We implement each round using standard data-parallel primitives
provided by \textsc{ParlayLib}~\cite{parlaylib}. The algorithm terminates
when no round discovers a new merge, exactly mirroring the fixed point of
the sequential procedure.

The two algorithms differ in how they detect new congruences. The algorithm
\parentalgo maintains a parent list for each term, and when a term is involved
in a union, it adds all of its parents to a worklist. The second algorithm
\filteralgo marks terms as dirty when they are involved in a union and then
filters all terms to find those whose children are dirty.

 
In summary, this paper makes the following contributions:

\begin{itemize}
    \item We present two parallel congruence closure algorithms, \parentalgo
    and \filteralgo, both built on a bulk-synchronous parallel design that
    replaces the classical sequential worklist with rounds of data-parallel
    primitives over a lock-free concurrent union-find.
    \item We characterize when our algorithms achieve good speedup using the
    notions of congruence \emph{depth} and \emph{width}, and show empirically
    that width is the dominant factor.
    \item Despite congruence closure being P-complete (i.e. believed to not
    have a theoretically efficient parallel
    algorithm)~\cite{unification}, we demonstrate near-linear speedups up to
    32 cores over an efficient
    sequential baseline on random benchmarks, a synthetic benchmark suite
    designed to stress-test our algorithm, and a set of circuit equivalence
    benchmarks~\cite{clausal-cc}.
    \item We open source our \href{code and benchmarks}{https://github.com/amarshah10/ParallelEgraph}.
\end{itemize}
% We additionally discuss what would be required to integrate our algorithms
% into an existing SMT solver or equality saturation engine.



\section{Background}
\label{sec:background}


\subsection{Congruence closure}

\emph{Congruence} ($\equiv$) requires that two terms that
that have the same function and pairwise-congruent children, must
be congruent, i.e. $f(s_1,
\ldots, s_k) \equiv f(t_1, \ldots, t_k)$ whenever $s_i \equiv t_i$ for every
$i$. \emph{Congruence closure} refers to the problem of building the smallest
equivalence relation that contains a given set of equalities and is closed under
congruence.

Typically, congruence closure is implemented using a
datastructure called an \emph{e-graph}~\cite{egraphs}, which maintains a
partition of terms into equivalence classes. More formally we consider:

\begin{definition}[]
    We define $T$ to be a set of terms over a list of function symbols $F$,
    where each $t\in T$ is either a leaf-term, for example $x, y, z$ or a
    compound term of the form $f(t_1, \ldots, t_k)$ where $f\in F$ $\text{arity}(f) = k$
    and $t_1, \ldots, t_k \in T$.
\end{definition}

We say that term $x$ is the \emph{child} of term $f(x)$ and conversely term
$f(x)$ is the \emph{parent} of term $x$.

\begin{definition}[Congruence Closure]
    Given $(F, T, E)$ where $F$ is a list of function symbols, $T$ is a set of
    terms over $F$, and $E = \{a_1 = b_1, \ldots, a_n = b_n\}$ is a set of
    equalities between terms in $T$ where $a_1, \ldots, a_n, b_1, \ldots, b_n
    \in T$ for every $i$, the congruence closure of $(F, T, E)$ is the smallest
    equivalence relation $\equiv$ over $T$ such that:
    \begin{enumerate}
        \item $a_i \equiv b_i$ for every $i$.
        \item $s \equiv s$ for every term $s$ (reflexivity).
        \item If $s \equiv t$ then $t \equiv s$ (symmetry).
        \item If $s \equiv t$ and $t \equiv u$ then $s \equiv u$ (transitivity).
        \item If $s_i \equiv t_i$ for every $i$ then $f(s_1, \ldots, s_k) \equiv
        f(t_1, \ldots, t_k)$ (congruence).     
    \end{enumerate}
\end{definition}

Congruence closure is useful in practice for SMT solvers, given a formula
$\phi = a_1 = b_1 \land \ldots \land a_n = b_n \land s \neq t$, an SMT solver
will run congruence closure on the equalities $E = \{a_i = b_i\}_{i=1}^n$ to
obtain the congruence closure $\equiv$ of $E$. Then, it will check if $s$ and
$t$ are in the same equivalence class under $\equiv$. If they are, then we 
know that $s \neq t$ is false and thus the formula $\phi$ is unsatisfiable.

Similarly, an equality saturation engine may ask: what programs are equivalent
to $P$ under a set of rewrite rules $E$. To answer this question, it can run
congruence closure on $E$ to obtain the congruence closure $\equiv$ of $E$.
Then, it may check which programs $Q$ are in the same equivalence class as $P$
under $\equiv$.

The classical algorithm of Nelson and Oppen~\cite{nelson-oppen-cc} maintains a
union-find over equivalence classes of subterms. A worklist holds classes that
have just been merged. Processing a class involves walking its \emph{parent
list}---all e-nodes that mention this class as a direct
child---re-canonicalizing each, and merging any pair that the hashcons
identifies as a duplicate. New merges produce new worklist entries; the loop
runs to a fixed point. An improvement by Downey et al.~\cite{downey} uses a 
hash-cons to detect new congruences in $O(1)$ time.

\subsection{Motivating Example}
\label{sec:bg:motivating}

To make the structure concrete, consider the formula which is solved by an SMT
solver:
%
\begin{equation*}\label{ex:formula}
\begin{aligned}
& \;\;\;\;\;\; a_0 = b_0 \land\; a_1 = b_1 \\
&\;\land\;g(f(a_0), f(a_1)) \neq g(f(b_0), f(b_1))
\end{aligned}
\end{equation*}
%

This formula consists of two equalities $a_0 = b_0$ and $a_1 = b_1$, and a
single disequality between two terms built from $f$ and $g$. The disequality
compares two applications of $g$, whose four arguments are the $f$-terms.

Congruence closure proceeds as follows. After applying the two asserted unions,
the union-find places $a_0 \equiv b_0$ and $a_1 \equiv b_1$. Then by congruence
we have that $f(a_0) \equiv f(b_0)$ and $f(a_1) \equiv f(b_1)$.

Once these equalities fire, the $g$-terms become congruent $g(f(a_0), f(a_1))
\equiv g(f(b_0), f(b_1))$. Thus, the SMT solver knows that these two terms must
be equal and thus the given formula is false.

Notice that this propagates in rounds: the initial equalities are added to the
union-find, then the four $f$-term equalities are added, and finally the
$g$-term equality is added.

\begin{definition}[Congruence Depth]
    We say that two terms are $k$-step congruent if they are congruent after $k$
    rounds of congruence closure. We say that the \emph{congruence depth} of a
    formula is the smallest $k$ such that every pair of congruent terms in the
    formula is $k$-step congruent.
\end{definition}

\begin{definition}[Congruence Width]
    For each round $r$ of congruence closure, we say that $\text{width}(r)$ is
    the number of new equalities that are added in round $r$. We define the
    total width of a congruence closure problem to be $\max_{r \in
    0\ldots\text{depth}} \text{width}(r)$, i.e.the maximum number of new
    equalities that can be added in a single round of congruence closure.
\end{definition}

In general, if we have a large number of terms that
become congruent in a single round, we can add all of these terms to the
union-find concurrently. This motivates our parallel algorithm for congruence
closure, which processes each round of congruence closure in parallel.


% : initially, the only equalities we know
% are $a_0 = b_0$ and $a_1 = b_1$, these can be added to the union-find
% concurrently. Once these are added, the four equalities between the $f$-terms
% are directly implied by these. Hence, they can be added concurrently in the next
% round. Finally, once these are added, the two $g$-terms are equal and can be
% added. Using our parallel algorithm (described in \autoref{sec:alg}), we can
% finish congruence closure in three rounds of size $2, 4, 1$ respectively. This
% is opposed to the sequential algorithm, which would process $7$ merges one at a
% time. We make these intuitions precise in Section~\ref{sec:alg}.


\subsection{Parallel algorithms}


\section{Algorithms}
\label{sec:algorithm}


\subsection{Sequential Baseline}~\label{sec:algorithm:no}

Our baseline is a simple yet efficient worklist algorithm based on one developed by Nieuwenhuis and Oliveras.

We maintain a \emph{signature table} that maps every term to the representative of its e-class. It is implemented via a hashtable, where hashing a term entails hashing its function symbol and the e-class identifiers of its operands. The signature table admits operations $\textsc{Call}(e)$ which enters term $e$ into the signature table, $\textsc{Lookup}(e)$, which finds the 



\subsection{Other Sequential Algorithms}~\label{sec:algorithm:topo}


\subsection{Bulk-Synchronous Parallel Algorithm}~\label{sec:algorithm:parallel}

\subsection{Asynchronous Parallel Algorithm}~\label{sec:algorithm:async}

\section{Implementation}
\label{sec:implementation}

\as{copied from final report -> not sure if we need to discuss these}

Here we describe some small implementation choices that had a large bearing on
performance. One was memoizing the hash function of a given term during semisort
to avoid recomputing it multiple times, which can be expensive. Hashing a term
requires hashing its function symbol, alongside the representatives of each of
its operands. This requires executing a \textsc{Find} for every term. We
therefore compute the hash of every term in the frontier at the beginning of a
round, store this in a \texttt{struct}, and perform semisort on these structs. 


Another optimization is that, when semisorting to determine congruence classes,
we perform an equality check on a secondary hash rather than deterministically
checking whether two terms are actually operand-wise equal as in the pseudocode
for \textsc{Congruent-Level1} in Figure~\ref{fig:alg:bsp}. This is because we
empirically found that having semisort probe over every bucket using this
deterministic equality check was very expensive. Changing the equality check to
a secondary hash allows the equality check when probing over a bucket to be much
cheaper. However, this introduces technical unsoundness, but this is tantamout
to a 96-bit hash collision, with probability about $10^{-14}$ for workloads of
our size (this can be computed via birthday-paradox like tricks). Although this
introduces a negligible chance of unsoundness, this change resulted in a 50\%
speedup, which we believe to be well-worth it.

\section{Evaluation}
\label{sec:evaluation}

\subsection{Random Benchmarks}

\subsection{Synthetic Benchmarks}

\subsection{Equality Saturation Benchmarks}

To get equality saturation benchmarks, we ran egglog~\cite{egglog} on the egg-cc
project~\cite{egg-cc}—an experimental optimizing compiler for the LLVM-like IR 
Bril. We ran egglog on a number of Bril programs for \as{todo} seconds and logged
all of the unions egglog generated during this time. We then extracted these unions
into and MST-LIB file and evaluated our parallel congruence closure algorithms on these
benchmarks.

\subsection{SMT Benchmarks}

To get SMT benchmarks, we took benchmarks from the QF-UF logic
in the annual 2025 SMT competition~\cite{smtcomp2025}. In this category, we picked
the benchmarks from \asm{fill in} as the other benchmarks were either very simple
or required more complicated boolean reasoning.

These benchmarks include boolean connectives, so to get pure congruence closure
constraints, we ran each of the bencmarks through the Sundance-SMT solver
~\cite{sundance}, every time Sundance-SMT found a conflict in its congruence 
closure solver, we extracted the set of equalities into an SMT-LIB file.







\section{Discussion}
\label{sec:discussion}

\subsection{Answer to Research Questions}~\label{sec:disc:rq}

We answer the research questions as follows:

\begin{itemize}
  \item \textbf{RQ1:} Our parallel congruence closure algorithm achieves a
  meaningful speedup on random benchmarks, a synthetic benchmark suite designed
  to stress test our algorithm, and a set of circuit equivalence benchmarks. We
  scale well up to 32 cores (\as{give a number here}), but began to plateau
  beyond that, especially on smaller benchmarks.
  \item \textbf{RQ2:} Our algorithm tends to perform better on benchmarks with a
  large number of merges per round. More rounds of congruence closure does not
  have a major effect on parallel speedup.
  \item \textbf{RQ3:} Our algorithm provides a meaningful speedup over a
  sequential implementation on the circuit equivalence benchmarks. It is
  difficult to evaluate our algorithm on SMT or equality saturation benchmarks,
  as congruence closure is part of a larger tool and thus difficult to isolate.
  We leave it to future work to integrate our algorithm into a state-of-the-art
  SMT solver or equality saturation engine.
\end{itemize}

\subsection{Algorithm Modifications and Future Work}~\label{sec:disc:limit-alg}

There are several extensions to the algorithm as presented to be addressed in 
future work.

\textbf{Synchronous Round Structure.} Our synchronous round structure could be
limited in terms of parallel scaling. In particular, if we have congruence
closure problems with large \emph{depth}, but small \emph{width}, our
parallization strategy cannot achieve a good speedup. We implemented an
asynchronous version of the dirty children algorithm, where we did semisort and
unions happened at the same time. Certain threads would semisort and discover
new congruences which the would add to a concurrent bag
datastructure~\cite{concurrentbag}. Then the threads that wer unioning would
probe this concurrent bag for new merges to perform.

We found that this was slower than the synchronous version. A challenge to
address is that even this asynchronous algorithm cannot efficiently parallize a
chain of congruences, i.e. where we you have add the merges $x \equiv y$, $f(x)
\equiv f(y)$, $f(f(x)) \equiv f(f(y))$, etc.



% inherently
% limits parallel scaling. In particular, the instance size processed at each
% round is small relative to the size of the problem for several benchmarks we
% have examined (like the \texttt{egg} benchmarks). We will likely have to take a
% more asynchronous approach to achieve better thread scaling; this would entail
% generating new congruences from discovered equalities immediately instead of
% waiting for the round to complete.

\textbf{Backtracking.} 
\as{Zak can you fill in how we could make our implementation backtrackable}

\textbf{Proof Production.} Our implementation does not produce \emph{proofs}(i.e.
\emph{witnesses}) of the equalities it discovers. For instance if our algorithm
discovers that $x = y$, it cannot replay the chain of equalities and congruences
that led to this conclusion. This is important for a number of reasons: it is
necessary to justify correctness and it is used to produce conflict clauses in
SMT solvers. Our implementation could easily be extended to support proof
production by avoiding path compression in the union-find data structure.
Another option is to maintain two union-find data structures: one for the actual
congruence closure algorithm, and another that maintains the proof forest as
done by the Z3 SMT solver does~\cite{z3-internals}.

\textbf{Incremental Update.} Our algorithms could support incremental update
(where equalities and new terms may arrive dynamically)

Finally, we would like to support our
implementation to incremental updates (where equalities and new E-node may
arrive dynamically). This is important for both equality saturation and SMT
solvers, where new equalities may be discovered via for example quantifier
instantiation.



\bibliographystyle{IEEEtran}
\bibliography{refs}

\end{document}